\chapter*{如何使用本仓库}
\addcontentsline{toc}{chapter}{如何使用本仓库}
\markboth{如何使用本仓库}{如何使用本仓库}

\section*{仓库结构}

\begin{lstlisting}[style=growthnoline, language=bash, caption={仓库目录结构}]
SGD/
 |-- main.tex                  # master file: cover, TOC, \input all chapters
 |-- notesstyle.sty            # shared style: layout, tcolorbox, listings, math
 |-- references.bib            # bibliography (bibtex + gbt7714 national standard)
 |-- Makefile                  # make / make check / make clean
 |-- tools/check.py            # repo linter: bad chars, cites, env pairing, counts
 |-- README.md                 # quick start
 |-- frontmatter/              # preface and how-to-use
 |-- chapters/                 # 21 chapters
 |-- appendix/                 # appendices A-E
 |-- code/                     # runnable example scripts
 `-- build/                    # build output (PDF)
\end{lstlisting}

\section*{编译方式}

\subsection*{方式一：Makefile（推荐）}

\begin{lstlisting}[style=growthnoline, language=bash]
make            # full build: xelatex -> bibtex -> xelatex -> xelatex
make quick      # single xelatex pass, for layout preview
make watch      # rebuild on save
make check      # repo lint: bad chars, cite keys, env pairing, word count
make clean      # remove intermediate files
make cleanall   # remove everything including the PDF
\end{lstlisting}

\subsection*{方式二：latexmk 手动}

\begin{lstlisting}[style=growthnoline, language=bash]
latexmk -xelatex -interaction=nonstopmode main.tex
\end{lstlisting}

\subsection*{方式三：Overleaf}

上传整个目录，把编译器切换为 \textbf{XeLaTeX}，并在项目设置里确认主文档为 \code{main.tex}。
若中文字体缺失，把 \code{main.tex} 中的 \code{fontset=macnew} 改为 \code{fontset=fandol}
（Overleaf 自带 Fandol 字体）。

\section*{三条学习路线}

\begin{table}[htbp]
\centering
\caption{按目标选择阅读顺序}
\label{tab:routes}
\small
\begin{tabularx}{\textwidth}{L{2.6cm} L{5.2cm} Y}
\toprule
\tabhead{路线} & \tabhead{章节组合} & \tabhead{适用人群与产出} \\
\midrule
算法工程师 &
\cref{ch:math} $\to$ \cref{ch:recall} $\to$ \cref{ch:ranking} $\to$ \cref{ch:sequence} $\to$ \cref{ch:receng}
$\to$ \cref{ch:search} $\to$ \cref{ch:rag} $\to$ \cref{ch:admodel} &
能独立完成「召回—粗排—精排—重排」链路设计与迭代，能讲清 CTR/CVR 建模与多目标融合的取舍。 \\
\addlinespace
增长操盘手 &
\cref{ch:overview} $\to$ \cref{ch:seo} $\to$ \cref{ch:geo} $\to$ \cref{ch:google} $\to$ \cref{ch:meta}
$\to$ \cref{ch:cro} $\to$ \cref{ch:growth} &
能设计自然流量 + 付费流量的组合打法，读懂后台指标，搭出从投放到回收的归因与实验闭环。 \\
\addlinespace
跨境卖家/供应链 &
\cref{ch:cbec} $\to$ \cref{ch:platform} $\to$ \cref{ch:supply} $\to$ \cref{ch:meta}
$\to$ \cref{ch:growth} &
能完成平台选择、选品验证、备货与履约方案设计，理解平台全托管/半托管的真实成本结构。 \\
\bottomrule
\end{tabularx}
\end{table}

\section*{全书地图}

\begin{figure}[htbp]
\centering
\begin{tikzpicture}[
  node distance=6mm,
  every node/.style={font=\sffamily\scriptsize, align=center},
  layer/.style={draw=#1, fill=white, rounded corners=2pt, minimum height=9mm, inner sep=4pt, text width=2.5cm},
  stage/.style={draw=NavyMid, fill=BoxBG, rounded corners=2pt, minimum height=9mm, inner sep=4pt, text width=2.9cm}
]
% 需求层
\node[stage] (demand) {需求产生\\搜索 / 推荐 / 广告位};
% 匹配层
\node[stage, below=8mm of demand] (match) {匹配与排序\\召回 · 粗排 · 精排 · 重排};
% 供给层
\node[stage, below=8mm of match] (supply) {供给与创意\\商品 · 内容 · 素材 · 落地页};
% 履约层
\node[stage, below=8mm of supply] (fulfill) {交易与履约\\支付 · 物流 · 关税 · 售后};
% 数据层
\node[stage, below=8mm of fulfill, fill=BoxBGAlt] (data) {数据与实验\\埋点 · 归因 · A/B · 因果};

\draw[-{Stealth}, thick, NavyMid] (demand) -- (match);
\draw[-{Stealth}, thick, NavyMid] (match) -- (supply);
\draw[-{Stealth}, thick, NavyDeep] (supply) -- (fulfill);
\draw[-{Stealth}, thick, TealDeep] (data.east) .. controls +(14mm,0) and +(14mm,0) .. (match.east);
\draw[-{Stealth}, thick, TealDeep] (data.west) .. controls +(-14mm,0) and +(-14mm,0) .. (match.west);
\node[right=16mm of demand, text width=3.1cm, draw=RuleGray, dashed, rounded corners=2pt, inner sep=4pt]
  (book) {本仓库覆盖：\\ \cref{ch:overview}--\cref{ch:growth}};
\draw[-{Stealth}, dashed, SlateGray] (book.west) -- (demand.east);
\end{tikzpicture}
\caption{「需求—匹配—供给—履约—数据」五层链路，本仓库的组织逻辑}
\label{fig:bookmap}
\end{figure}

\section*{写作与贡献约定}

\begin{itemize}
  \item 每章以 \code{chapters/chNN-*.tex} 为独立文件，只写正文，不写导言区；
  \item 提示框统一使用 \code{keypoint / insight / casebox / pitfall / toolbox / interview / formulacard} 环境；
  \item 引用统一走 \code{references.bib}，新增文献请补 key 后再引用；
  \item 代码块内\textbf{只使用 ASCII 字符}（中文注释请写在正文），避免 CJK 在 \code{listings} 中渲染异常；
  \item 表格宽度优先用 \code{tabularx} + \code{Y} 列，避免溢出版心。
\end{itemize}
